\begin{abcliste}
\abc The power is the intensity times the area:
\begin{align}
 P &= \PeF = \Ia\times\left(\da\right)^2 = \Pe \approx \result{\PenP{0}{2}\cdot 10^{-23}\,\si{\watt}}
\end{align}

\abc The electron needs at least the work function $W = \WaO = \Wa$:
\begin{align}
 t &= \twF = \frac{\Wa}{\Pe} = \tw \approx \result{\twP{0}{2}}
\end{align}
That is about $\SI{2.5}{\hour}$. In the wave model, dim light would release the first electrons only after hours. But they leave the plate $\result{\text{at once}}$ (within less than \SI{e-9}{\s}): the wave model fails.

\abc The light delivers its energy in portions, photons, each with
\begin{align}
 E_\mathrm{ph} &= \EpVF = \frac{\nch\times\ncc}{\nce\times\la} = \EpV \approx \result{\EpVP{0}{3}}
\end{align}
which is more than the work function of zinc. Dim light means few photons, not weak ones. On \AaO\ there arrive
\begin{align}
 N &= \NrF = \frac{\Ia\times\Aa}{\Ep} = \Nr \approx \result{\NrnP{0}{2}\cdot 10^{9}\,\si{\per\second}}
\end{align}
photons per second. The first photon to hit an electron gives it its whole energy at once and can free it immediately; nothing has to be collected.
\end{abcliste}
